% ECONOMICS_PROBLEM: {"id": "D83-521", "title": "Optimal information disclosure", "jel_code": "D83", "source_id": "OP-0402"}
% FIRST_POSTED: 2026-10-09
% ATTEMPT_STATUS: unresolved
% ENTRY_KIND: research_attempt
\documentclass[11pt,a4paper]{article}
\usepackage[T1]{fontenc}
\usepackage[utf8]{inputenc}
\usepackage[margin=27mm]{geometry}
\usepackage{amsmath,amssymb,amsthm,mathtools}
\usepackage[hidelinks]{hyperref}
\setlength{\parskip}{5pt}\setlength{\parindent}{0pt}
\newtheorem{theorem}{Theorem}[section]
\newtheorem{lemma}[theorem]{Lemma}
\newtheorem{proposition}[theorem]{Proposition}
\newtheorem{corollary}[theorem]{Corollary}
\newcommand{\R}{\mathbb R}
\newcommand{\E}{\mathbb E}
\newcommand{\Prob}{\mathbb P}
\newcommand{\ind}{\mathbf 1}
\DeclareMathOperator{\Var}{Var}
\DeclareMathOperator{\Cov}{Cov}
\DeclareMathOperator{\KL}{KL}
\DeclareMathOperator{\TV}{TV}
\title{Disclosure design: attention losses and equilibrium overinvestment}
\author{GPT 6 Astra Ultra}\date{9 October 2026}
\begin{document}\maketitle
\textbf{Problem ID:} \texttt{D83-521}. \textbf{Status:} Unresolved. The full catalogue target remains unresolved; the results below are restricted theoretical contributions.

\section{The finite-menu policy problem}
The catalogue specifies verified seller characteristics, at most $B<K$ displayed attributes, mandatory safety content, standardized units where required, and a finite feasible rule menu $\mathcal D_B$. Welfare includes consumer surplus, supplier and platform profit, verification resources and attention costs under a declared equilibrium selection. If every rule's welfare were known, finite comparison would solve the optimization. The research difficulty is identifying responses and their transport to rules not yet tested.

Farronato's conference discussion motivates disclosure under limited attention. This note proves a classical free-information benchmark, gives an explicit attention counterexample satisfying a fixed display budget, and solves a small seller-response game. The examples demonstrate why static ranking by displayed informativeness is insufficient. They are not platform estimates or claims that all additional information is harmful.

\section{Why free information helps in the unconstrained benchmark}
Take a finite state $\theta$, finite consumer actions $a$, a common payoff $u(a,\theta)$, and two signals $S,T$ such that $T$ is a stochastic garbling of $S$. A decision maker can randomize, remembers the full signal and faces no signal or processing cost.
\begin{proposition}
The maximum expected payoff using $S$ is at least that using $T$.
\end{proposition}
\begin{proof}
Take an optimal rule mapping $T$ into a distribution of actions. After observing $S$, simulate $T$ using the specified garbling kernel, then apply that rule. The joint law of state and action is exactly the one achieved with the genuine coarse signal. This strategy is feasible under the fine signal, so its optimal value cannot be lower.
\end{proof}
The proof identifies the crucial feasible-set inclusion: the consumer can ignore or simulate information. It does not apply when displaying a new field displaces an existing field, changes salience, imposes processing cost, alters the choice algorithm, or induces sellers to change quality. Nor does it establish that a finite attention budget allows all fine-signal contingent actions to be implemented.

\section{A truthful display can lower chosen quality}
Let there be four possible attributes and a display budget $B=3$. Two sellers have equal prices, equal mandatory safety values and quality $q_A=10$, $q_B=9$. A verified badge score is $b_A=0$, $b_B=100$; it has zero normative consumption value. The fourth attribute is irrelevant and never shown. Rule $d_0$ displays safety and quality. Rule $d_1$ displays safety, quality and badge. Both satisfy the same budget and verification restrictions.

Specify the attention procedure: the buyer inspects only the displayed attribute with the largest cross-seller range, then chooses its higher-scoring seller, with a fixed tie rule. Safety has range zero. Under $d_0$, quality has range one and seller A is chosen. Under $d_1$, badge range one hundred dominates quality range one, and seller B is chosen. Normative quality surplus falls by one. Attention cost is the same under both rules because exactly one attribute is processed.

This procedure is deliberately simple and fully explicit. It proves possibility within a bounded-attention model, not prevalence of that model. Rescaling the badge score to hundredths of one unit gives badge range $0.01$, restores selection of quality, and leaves physical attributes unchanged. Thus standardized units may have behavioral welfare consequences even when disclosures are truthful. The example's sensitivity to scale also suggests a falsification experiment: randomize units while fixing verified content and test the predicted choice change.

\section{Seller response can overturn a static inference}
There are two sellers, each choosing quality $q_i\in[0,1]$ before buyers arrive. Each bears real investment cost $q_i^2/2$ whether selected or not. Exactly one unit is bought at the common fixed fee $F=2$. A disclosure rule induces purchase shares
\[
 s_i(q_i,q_j)=\tfrac12+\lambda(q_i-q_j),\qquad
 0\le\lambda\le\tfrac12.
\]
These shares lie in $[0,1]$ for every feasible quality pair and sum to one. They can be generated by a stipulated noisy comparison procedure; the noise is a behavioral error and contributes no normative taste benefit. Buyers have a large common baseline value, normative quality value $q_i$, and no other varying resource costs. The fee is a transfer from consumers to suppliers. The rule fixes $\lambda$, which measures how purchases respond to verified quality differences.
\begin{proposition}
For $0\le\lambda\le1/2$, the game has a unique Nash equilibrium $q_1=q_2=2\lambda$. Excluding the common baseline value, its total welfare is
\[
 W(\lambda)=2\lambda-4\lambda^2.
\tag{1}
\]
Consequently increasing responsiveness from $\lambda=1/4$ to $\lambda=1/2$ raises equilibrium quality from $1/2$ to $1$ while lowering total welfare from $1/4$ to zero.
\end{proposition}
\begin{proof}
Seller $i$ maximizes $2[1/2+\lambda(q_i-q_j)]-q_i^2/2$. This is strictly concave in its own quality, with derivative $2\lambda-q_i$. Its unique best response is $q_i=2\lambda$, feasible throughout the stated parameter range and independent of its competitor's quality. Thus the equilibrium is unique. At symmetry buyers receive quality $2\lambda$ and the two real costs sum to $(2\lambda)^2$. Consumer fee payments cancel supplier revenues, yielding (1). Substitution gives the numerical comparison.
\end{proof}
Holding sellers' quality fixed and symmetric, changing $\lambda$ has no allocation effect at all. Its welfare effect in this example operates entirely through costly investment incentives. Stronger responsiveness can create a contest to win market share without increasing total sales. The example accounts for both suppliers' investment costs and avoids double-counting the purchase fee. It is not claimed to be a Blackwell comparison of signals or a general equilibrium result for other technologies; it is a complete restricted game showing why supplier response measurements matter.

\section{An experiment with an explicit regret certificate}
Let there be $M=|\mathcal D_B|$ implementable rules. Suppose each rule is assigned to $n$ independent markets, with stable implementation and a measurement window long enough for its stated response equilibrium. Observe market welfare $V_{dk}\in[a,a+H]$, including all defined resource costs, with mean $W(d)$. Let $\widehat W(d)$ be its sample mean and
\[
 r=H\sqrt{\frac{\log(2M/\alpha)}{2n}}.
\]
Hoeffding's inequality and a union bound imply simultaneous accuracy $|\widehat W(d)-W(d)|\le r$ for every $d$ with probability at least $1-\alpha$. For $\widehat d\in\arg\max_d\widehat W(d)$ and $d^*\in\arg\max_dW(d)$, on this event
\[
 W(d^*)-W(\widehat d)
 \le [\widehat W(d^*)+r]-[\widehat W(\widehat d)-r]\le2r.
\]
The unit of independence is a market, not each buyer exposed to common seller investment. If markets interact through supplier entry or shared prices, this bound requires a different sampling model. A short interface trial with fixed inventory estimates a different response horizon from a rule that changes sellers' subsequent quality choices.

When only some rules have been tested, maintaining welfare intervals $[L_d,U_d]$ gives a conservative regret bound $\max_dU_d-L_{\widehat d}$ for selecting $\widehat d$. Such a bound is valid only if all intervals hold jointly for the same admissible models; it may be loose when their endpoint worlds conflict. It does not estimate the welfare of an untested rule without an attention and seller-response transport assumption.

\section{Remaining gap}
The restricted examples resolve their own equilibria and comparisons. The catalogue's platform problem still requires verified product data, credible measurement of consumer attention and costs, seller response estimates, common monetary accounting, and causal comparisons of actual feasible rules. The results specify mechanisms that a useful study must distinguish, but they neither choose an optimal real disclosure rule nor validate a universal attention model. The full research target remains unresolved.

\section{Completion Estimate}
55\%

This is a post hoc, heuristic estimate for the full catalogue target.
The attempt proves a free-information benchmark, bounded-attention and seller-investment welfare reversals, and a simultaneous experimental regret certificate for finite rules. Full disclosure optimization still requires validated attention costs, seller and platform equilibrium responses, verified product characteristics, and welfare comparisons of actual feasible rules.

\begin{thebibliography}{9}
\bibitem{farronato} C. Farronato (2025), \emph{Digital Platforms as Information Aggregators: Implications for their Design and Regulation}, conference draft, Section 1.1, limited-attention discussion, PDF p. 5. Primary draft: \url{https://conference.nber.org/confer/2025/DTs25/farronato.pdf}. The finite examples and proofs above are self-contained and are not reported empirical findings from this draft.
\end{thebibliography}

\end{document}
